\documentclass[12pt,a4paper]{report}
\usepackage[utf8]{inputenc}
\usepackage[T1]{fontenc}
\usepackage[provide=*,french]{babel}
\usepackage{geometry}
\usepackage{hyperref}
\usepackage{graphicx}
\usepackage{array}
\usepackage{booktabs}
\usepackage{longtable}
\usepackage{caption}
\usepackage{setspace}

\geometry{left=2.5cm,right=2.5cm,top=2.5cm,bottom=2.5cm}
\setstretch{1.15}
\hypersetup{colorlinks=true,linkcolor=black,urlcolor=blue}

\begin{document}

\begin{titlepage}
\centering
{\Large Royaume du Maroc \par}
\vspace{0.8cm}
{\Large Projet académique de génie logiciel et cybersécurité \par}
\vspace{2.0cm}
{\Huge \textbf{ChabahRoot} \par}
\vspace{0.5cm}
{\LARGE Plateforme de télémétrie noyau Linux fondée sur tracefs \par}
\vspace{1.3cm}
{\large Rapport détaillé de consolidation architecturale \par}
\vspace{0.8cm}
{\large Date de rédaction : 13 mai 2026 \par}
\vfill
{\large Auteur institutionnel : projet ChabahRoot \par}
\vspace{0.3cm}
{\large Version étudiée : arborescence finale francisée \par}
\vfill
{\large \today \par}
\end{titlepage}

\pagenumbering{roman}

\chapter*{Résumé}
\addcontentsline{toc}{chapter}{Résumé}
Le présent rapport étudie ChabahRoot comme une plateforme de télémétrie noyau Linux construite autour de tracefs, dans un contexte où le dépôt se trouvait dans un état intermédiaire mêlant anciennes conventions de nommage, répertoires dupliqués, scripts partiellement réécrits et résidus d’une ancienne direction eBPF. L’objectif du travail n’a pas été d’ajouter une nouvelle fonctionnalité spectaculaire, mais de restaurer une cohérence globale entre l’architecture du projet, les chemins réellement exécutables, les points d’entrée shell et la documentation.

La version finale s’organise autour de cinq modules en français. Veille\_kernel concentre l’interaction avec le noyau et l’activation des tracepoints. Chaines\_d\_ecoute fournit un schéma d’événement stable et joue le rôle de normalisateur. Vigie\_comportementale porte la lecture comportementale à travers un moteur de règles et deux mécanismes complémentaires de contextualisation. Chambre\_d\_analyse relie les étages opérationnels et offre une entrée unique pour l’exploitation hors ligne ou l’intégration à une source active. Enfin, socle\_commun rassemble les mécanismes de journalisation et de configuration transverses.

La thèse défendue est qu’une architecture volontairement sobre peut devenir plus utile qu’une architecture plus ambitieuse mais inachevée. Le nettoyage final a montré que, dans un projet shell proche du système, la qualité des chemins, des frontières de responsabilité et de la séparation entre messages humains et données de traitement pèse souvent plus lourd que la sophistication théorique des composants. ChabahRoot devient ainsi un cas d’étude sur la consolidation d’un pipeline système crédible, lisible et maintenable.

\chapter*{Abstract}
\addcontentsline{toc}{chapter}{Abstract}
This report analyzes ChabahRoot as a Linux kernel telemetry platform based on tracefs. The repository had been left in a broken transitional state with mixed naming conventions, duplicate directories, obsolete eBPF remnants, and inconsistent shell entry points. The purpose of the final pass was not to expand functionality but to restore architectural truthfulness across the codebase, the runtime paths, and the documentation.

The final structure revolves around five French named modules. Veille\_kernel handles tracefs interaction and tracepoint activation. Chaines\_d\_ecoute normalizes raw observations into a stable event schema. Vigie\_comportementale applies behavioral rules and supporting host level checks. Chambre\_d\_analyse orchestrates the usable analysis path for offline and active inputs. Socle\_commun centralizes logging and shared configuration. The report argues that, for a shell based systems project, coherence, explicit interfaces, and reliable execution paths matter more than keeping half completed low level ambitions alive.

\tableofcontents
\listoffigures
\listoftables

\chapter*{Liste des acronymes}
\addcontentsline{toc}{chapter}{Liste des acronymes}
\begin{longtable}{>{\raggedright\arraybackslash}p{3cm} >{\raggedright\arraybackslash}p{10cm}}
API & Interface de programmation applicative. \\
NDJSON & Format dans lequel chaque ligne contient un objet JSON autonome. \\
PID & Identifiant numérique d’un processus sous Linux. \\
PPID & Identifiant numérique du processus parent. \\
SIEM & Plateforme de centralisation, de corrélation et d’analyse d’événements de sécurité. \\
SUID & Mécanisme Unix permettant l’exécution d’un binaire avec un identifiant effectif élevé. \\
UID & Identifiant numérique de l’utilisateur effectif. \\
\end{longtable}

\chapter*{Introduction générale}
\addcontentsline{toc}{chapter}{Introduction générale}
Le besoin de télémétrie noyau naît d’une tension bien connue dans les systèmes de sécurité. Plus l’observation se situe près du noyau, plus elle permet de voir des signaux robustes d’activité, comme des transitions d’identité, des changements de contexte d’exécution ou des tentatives d’élévation de privilèges. Dans le même temps, plus cette observation devient spécialisée, plus elle risque de dépendre d’outils complexes, de prérequis stricts et d’une maintenance exigeante. ChabahRoot se place précisément dans cet entre-deux. Le projet cherche à produire une surveillance noyau intelligible et directement exploitable, sans assumer le coût complet d’une instrumentation bas niveau sophistiquée.

Cette orientation explique le choix de tracefs. Le recours à une interface déjà intégrée au noyau et exposée sous forme de système de fichiers présente plusieurs avantages dans le cadre étudié. D’abord, la surface de déploiement reste faible. Ensuite, la lecture de l’état du système est plus simple, car les mécanismes d’activation sont visibles à travers des chemins explicites. Enfin, l’architecture shell du projet reste cohérente avec ce mode d’interaction. Le dépôt n’a pas à maintenir des outils supplémentaires pour compiler ou attacher un programme spécialisé au noyau ; il peut concentrer ses efforts sur la chaîne de traitement et sur la stabilité des interfaces entre modules.

Le rapport n’examine donc pas seulement un produit fini. Il examine une opération de consolidation. Le dépôt initial laissait apparaître un projet déjà partiellement réorienté vers tracefs, mais encore encombré par des traces d’une ancienne approche, par des renommages incomplets et par une coexistence confuse entre conventions anglaises et françaises. Une telle situation n’est pas anecdotique. Dans un projet shell proche du système, l’incohérence structurelle devient elle-même une source de panne. Les chapitres qui suivent montrent comment cette incohérence a été ramenée à une architecture lisible, quelles justifications ont guidé les choix retenus, quelles remarques doivent accompagner l’usage du système, et dans quelle mesure le résultat obtenu répond désormais à une logique technique unifiée.

\chapter{Fondements théoriques et justification du choix tracefs}
L’infrastructure Linux de traçage offre plusieurs voies d’observation. Certaines approches, comme eBPF, permettent une expressivité importante et une proximité très forte avec le noyau. D’autres approches privilégient la disponibilité et la régularité. Tracefs relève de cette seconde famille. Il donne accès à des points de trace prédéfinis sans imposer une chaîne de compilation dépendante du noyau exact, de ses en-têtes ou de son environnement de construction. Dans le contexte de ChabahRoot, cette propriété pèse lourd. Le projet vise une solution de télémétrie directement mobilisable, compréhensible à la lecture et suffisante pour démontrer une chaîne cohérente d’acquisition, de normalisation et de détection.

Ce choix doit être compris comme une décision d’architecture et non comme un renoncement abstrait. La branche eBPF qui existait dans le dépôt n’était plus simplement incomplète ; elle créait une fiction de double direction technique. En la maintenant, le projet se condamnait à transporter deux récits incompatibles : celui d’un pipeline shell orienté tracefs, et celui d’une instrumentation in-kernel plus ambitieuse mais non stabilisée. La justification du basculement définitif vers tracefs tient donc autant à la faisabilité qu’à la vérité du système. Un composant non aligné avec le chemin d’exécution réel devient une dette, même lorsqu’il semble prometteur sur le papier.

Du point de vue de la détection, les événements retenus possèdent une valeur sémantique forte. Execve marque un changement de programme qui peut révéler une exécution depuis un emplacement temporaire ou inattendu. Setuid et setgid exposent des transitions de privilèges ou de contexte d’identité. Prctl, dans certains usages, fournit une fenêtre sur des modifications de posture plus discrètes. Ces événements ne constituent pas une couverture exhaustive du comportement noyau, mais ils sont assez parlants pour soutenir un pipeline démonstratif et pour motiver des règles comportementales crédibles.

\begin{figure}[h]
\centering
\fbox{\parbox{0.82\textwidth}{\centering
L’architecture conceptuelle de ChabahRoot peut être lue comme un passage progressif d’un signal brut noyau vers une interprétation humaine. La capture tracefs occupe le point le plus bas, la normalisation impose une forme stable, la vigie comportementale transforme cette forme en signal de sécurité, puis la chambre d’analyse organise l’usage du tout dans une trajectoire opérationnelle unique.}}
\caption{Lecture conceptuelle de la chaîne de traitement ChabahRoot}
\end{figure}

\chapter{Exigences, contraintes et remarques de cadrage}
Les exigences fonctionnelles du projet apparaissent simples, mais elles s’enracinent dans des contraintes fortes. La première est la nécessité de produire un schéma d’événement stable quelles que soient les conditions d’entrée. Un pipeline qui manipule alternativement des corpus hors ligne, des flux shell et des traces système ne peut pas s’appuyer sur des structures implicites. La normalisation n’est donc pas un confort ; elle est le point de départ de la crédibilité du reste du système. Sans elle, chaque module devrait réinterpréter à sa manière l’origine et la signification des données reçues.

La seconde contrainte est d’ordre opérationnel. ChabahRoot doit pouvoir être essayé rapidement sans imposer une préparation lourde à chaque usage. Cette exigence conduit à distinguer les modes d’exploitation. La capture réelle dépend de root et de tracefs. En revanche, l’analyse hors ligne doit rester praticable à partir d’un corpus d’exemple. Cette distinction n’est pas un compromis défensif. Elle est une manière de protéger la vitesse d’itération. Sans cette possibilité de test hors ligne, toute évolution du pipeline serait ralentie par la dépendance à un environnement noyau exact.

Une troisième contrainte concerne la lisibilité du dépôt lui-même. Le projet ne pouvait pas être considéré comme stabilisé tant qu’il mêlait plusieurs arborescences concurrentes, plusieurs points d’entrée historiques et plusieurs styles de scripts. Dans un projet shell, la qualité de l’architecture se reflète immédiatement dans la facilité de lecture, dans la confiance que l’on accorde aux chemins, et dans la capacité à diagnostiquer un incident sans d’abord reconstituer l’intention cachée d’un renommage incomplet. La consolidation a donc visé à rendre les choix visibles et continus dans l’ensemble de l’arbre.

\begin{table}[h]
\centering
\begin{tabular}{>{\raggedright\arraybackslash}p{4cm} >{\raggedright\arraybackslash}p{9cm}}
\toprule
Contrainte & Effet concret sur le design final \\
\midrule
Disponibilité d’un point d’observation noyau simple & Abandon définitif de la chaîne eBPF au profit de tracefs \\
Besoin d’essais rapides hors environnement privilégié & Maintien d’un corpus d’exemple traversable sans accès root \\
Nécessité d’une détection lisible & Externalisation des règles dans un fichier JSON interprété par jq \\
Fragilité structurelle d’un pipeline shell & Discipline stricte sur les chemins, la journalisation et les points d’entrée \\
Volonté de cohérence documentaire & Alignement complet entre noms de modules, scripts exécutables et textes d’accompagnement \\
\bottomrule
\end{tabular}
\caption{Contraintes de cadrage et effets observables sur la solution retenue}
\end{table}

\chapter{Architecture finale et justification des modules}
L’architecture finale ne se contente pas de juxtaposer cinq répertoires ; elle distribue clairement les responsabilités. Veille\_kernel se tient au plus près du noyau. Ce module vérifie l’environnement, active les tracepoints visés, maintient le cycle de vie de la capture et fournit un flux brut à l’aval. Sa justification est d’ordre défensif : tout ce qui écrit dans tracefs doit être contenu dans un périmètre réduit et facilement auditable. Multiplier ces écritures à travers plusieurs scripts transverses aurait contredit le principe de séparation retenu.

Chaines\_d\_ecoute répond à une autre nécessité. Le signal brut du noyau n’est pas, en l’état, une matière satisfaisante pour la détection. Les lignes de trace doivent être interprétées, reformatées, enrichies et stabilisées. Ce module joue donc le rôle d’interface canonique entre le monde bas niveau et le monde analytique. Sa présence se justifie moins par la transformation de données en elle-même que par la suppression de toute ambiguïté entre la capture et la lecture métier. Le reste du système ne traite plus des traces ; il traite des événements.

Vigie\_comportementale constitue le cœur interprétatif du projet. Le moteur de règles y occupe la place centrale parce qu’il rend les critères de détection visibles et modifiables sans reconfigurer le contrôle shell. Ce choix prolonge la philosophie générale du projet : déplacer la variabilité dans des formes déclaratives chaque fois que possible. Les composants d’audit offensif et de surveillance défensive n’ont pas la même valeur probante que les événements noyau, mais ils apportent un contexte utile. Leur rôle est de compléter la lecture du système, non de brouiller la hiérarchie entre preuve de capture et indice d’exposition.

Chambre\_d\_analyse a une justification particulièrement importante dans la version consolidée. Le dépôt intermédiaire souffrait de chemins historiques morts et d’une dispersion des points d’entrée. En recentrant l’exécution autour d’un chemin d’analyse unique, ce module simplifie l’usage et supprime l’incertitude sur la bonne manière de traverser le pipeline. Enfin, socle\_commun fournit un service de cohérence. La journalisation, la configuration et la résolution de certains chemins critiques y sont regroupées pour éviter les dérives locales. Cette mutualisation n’est pas seulement pratique ; elle permet de limiter les écarts silencieux entre scripts qui, dans un environnement shell, finissent par se transformer en comportements contradictoires.

\begin{table}[h]
\centering
\begin{tabular}{>{\raggedright\arraybackslash}p{4cm} >{\raggedright\arraybackslash}p{3.6cm} >{\raggedright\arraybackslash}p{5.7cm}}
\toprule
Module & Fonction dominante & Justification principale \\
\midrule
veille\_kernel & Capture et gestion tracefs & Réduire la surface d’écriture sur le noyau à un périmètre lisible \\
chaines\_d\_ecoute & Normalisation NDJSON & Fournir une forme stable et indépendante de l’origine du flux \\
vigie\_comportementale & Détection et contexte & Séparer les critères de sécurité de la mécanique shell \\
chambre\_d\_analyse & Orchestration d’usage & Supprimer la dispersion des points d’entrée et unifier le parcours \\
socle\_commun & Journalisation et configuration & Empêcher les divergences locales dans les comportements partagés \\
\bottomrule
\end{tabular}
\caption{Rôle et justification de chaque module dans l’architecture finale}
\end{table}

\chapter{Décisions d’implémentation, nettoyage et remarques critiques}
Le nettoyage du dépôt a d’abord consisté à supprimer ce qui compromettait la vérité architecturale du système. Les répertoires dupliqués, les scripts artificiels, les chemins résiduels vers des structures abandonnées et la branche eBPF obsolète n’étaient pas seulement superflus ; ils introduisaient un risque permanent de mauvaise lecture et de mauvais lancement. En retirant ces éléments, le projet perd une illusion de richesse mais gagne une identité technique nette.

Une décision notable a concerné les commentaires shell. Dans un dépôt comme ChabahRoot, le commentaire n’est utile que s’il signale une intention ou une contrainte qui ne se voit pas immédiatement dans le code. Des commentaires génériques, décoratifs ou formulés comme un pseudo rapport automatique dégradent la lisibilité plus qu’ils ne l’améliorent. La réécriture en français ASCII a donc eu une portée plus large qu’une simple contrainte de style. Elle a contribué à rendre les scripts compatibles avec l’image d’un dépôt maintenu par des personnes attentives à la précision et non par une succession de couches générées.

La question des flux de sortie a également mérité une attention particulière. Dans une architecture shell, un log sur le mauvais flux peut corrompre une substitution de commande, un pipe ou une consommation aval. C’est pourquoi la séparation entre messages humains et données machine a été défendue comme un principe de premier ordre. Cette remarque vaut autant pour la capture noyau que pour les modules de plus haut niveau. L’intégrité d’un pipeline shell ne tient pas seulement à la syntaxe correcte de ses scripts, mais à la discipline invisible de ses flux.

Enfin, la résolution des chemins de journaux a révélé une difficulté typique des dépôts en transition. Un chemin relatif qui semblait convenable dans un ancien mode d’exécution pouvait devenir source d’échec dès qu’un autre point d’entrée lançait le même module depuis un répertoire différent. Le socle commun a donc été ajusté pour rendre ces chemins robustes et pour prévoir un repli exécutable lorsque le chemin principal n’est pas accessible. Cette correction est représentative de la nature du nettoyage final : moins orientée vers l’ajout de fonctionnalités que vers l’élimination de causes silencieuses d’instabilité.

\chapter{Expérimentation, résultats et interprétation}
L’expérimentation menée sur la version consolidée du dépôt a d’abord eu pour but de vérifier que les modules pouvaient être sollicités dans des conditions cohérentes avec leur rôle final. Le passage du corpus d’exemple dans le normaliseur a confirmé que la chaîne de transformation produit un schéma stable, enrichi et suffisamment uniforme pour alimenter la couche de détection. Cette observation est importante parce qu’elle valide la promesse centrale de chaines\_d\_ecoute : faire disparaître la variabilité d’origine au profit d’un événement exploitable.

La couche de détection a ensuite montré qu’elle pouvait charger les règles finales et réagir à plusieurs formes d’événements significatifs sur un même parcours. Les alertes observées sur les scénarios d’exemple correspondent à des catégories cohérentes avec les objectifs du projet : exécution depuis un emplacement temporaire, transition de privilège, modification de capacité et motif de commande jugé suspect. L’intérêt principal n’est pas dans la seule présence de ces alertes, mais dans le fait qu’elles émergent d’un schéma de données devenu suffisamment stable pour supporter une lecture déclarative continue.

L’orchestration par chambre\_d\_analyse a fourni un troisième résultat significatif. Elle prouve que l’usage du projet ne dépend plus de chemins historiques morts ni de services annexes devenus fictifs. Cette clarification améliore directement la qualité d’exploitation. Un utilisateur n’a plus à reconstruire la généalogie du dépôt pour savoir quelle commande déclenche la chaîne pertinente. En matière d’outillage système, cette simplicité apparente constitue une victoire concrète, car elle réduit le risque d’écart entre l’intention de l’opérateur et la réalité du chemin exécuté.

Le contrôle de veille\_kernel sans privilège root a enfin servi de test d’hygiène. Le fait que cette couche échoue rapidement et clairement lorsque le contexte n’est pas compatible avec une activation tracefs montre que le projet préfère l’explicitation au fonctionnement partiel. Cette propriété mérite d’être soulignée comme une remarque de sûreté. Dans les outils proches du système, un échec franc vaut souvent mieux qu’une réussite ambiguë qui laisserait croire à une capture active alors que les tracepoints ne sont pas réellement en service.

\begin{table}[h]
\centering
\begin{tabular}{>{\raggedright\arraybackslash}p{4cm} >{\raggedright\arraybackslash}p{4cm} >{\raggedright\arraybackslash}p{5cm}}
\toprule
Observation & Module concerné & Interprétation \\
\midrule
Stabilité du schéma d’événement & chaines\_d\_ecoute & Le normaliseur remplit correctement son rôle de frontière de forme \\
Alertes cohérentes sur corpus d’essai & vigie\_comportementale & Les règles reposent sur une structure devenue suffisamment robuste \\
Parcours d’usage unifié & chambre\_d\_analyse & Le dépôt ne dépend plus de points d’entrée historiques cassés \\
Échec explicite sans privilège root & veille\_kernel & La couche noyau protège l’opérateur contre un faux sentiment d’activation \\
\bottomrule
\end{tabular}
\caption{Lecture synthétique des résultats observés sur la version consolidée}
\end{table}

\chapter*{Discussion générale}
\addcontentsline{toc}{chapter}{Discussion générale}
Le principal enseignement de cette consolidation tient au rapport entre modestie technique et fiabilité perçue. Un projet peut paraître plus ambitieux lorsqu’il accumule plusieurs directions, plusieurs répertoires et plusieurs styles d’implémentation. Pourtant, cette apparente richesse se retourne rapidement contre lui lorsque rien ne garantit que ces éléments convergent vers un même usage réel. ChabahRoot montre qu’une réduction disciplinée du bruit structurel peut produire plus de valeur opérationnelle que le maintien d’éléments inachevés supposés prometteurs.

Il convient toutefois de formuler une remarque d’équilibre. Le choix tracefs ne doit pas être lu comme une supériorité absolue sur toute autre solution. Il doit être compris comme une réponse appropriée à la nature du projet. Pour un dépôt orienté shell, conçu comme plateforme de télémétrie lisible et déployable avec peu de friction, la cohérence l’emporte sur la sophistication. Dans un autre contexte, centré sur la performance extrême ou sur l’observation noyau approfondie, la conclusion pourrait être différente. Le mérite de ChabahRoot est d’avoir cessé de prétendre simultanément à deux modèles incompatibles.

Une autre discussion porte sur la documentation elle-même. Dans un dépôt en transition, la documentation a tendance à devenir soit trop pauvre, soit mensongère par obsolescence. Le travail accompli ici cherche à éviter ces deux défauts. Le README assume un rôle d’usage direct, tandis que le rapport porte une fonction plus analytique, destinée à exposer les justifications, les choix et les remarques critiques. Cette séparation des registres documentaires participe elle aussi à la cohérence du système.

\chapter*{Conclusion}
\addcontentsline{toc}{chapter}{Conclusion}
La consolidation finale de ChabahRoot a transformé un dépôt hybride en une plateforme de télémétrie noyau cohérente, lisible et alignée sur un unique chemin d’exécution. Le projet repose désormais sur une capture tracefs assumée, sur une normalisation clairement identifiable, sur une couche de détection fondée sur des règles visibles, sur une orchestration d’usage recentrée et sur un socle commun plus robuste. Cette cohérence rend le système plus crédible, plus maniable et plus honnête vis-à-vis de ce qu’il promet réellement.

La contribution principale de ce travail n’est donc pas l’ajout d’une capacité entièrement nouvelle, mais la restauration d’une continuité entre intention architecturale, implémentation, documentation et usage. Dans un domaine où les outils proches du système sont facilement fragilisés par des détails de structure, cette continuité constitue déjà une valeur substantielle. ChabahRoot devient ainsi une base saine pour de futurs enrichissements, précisément parce qu’il ne porte plus en lui plusieurs histoires techniques concurrentes.

\chapter*{Perspectives}
\addcontentsline{toc}{chapter}{Perspectives}
Les perspectives les plus naturelles concernent d’abord la sortie analytique du système. Le répertoire salle\_des\_rapports peut accueillir une production plus aboutie, à condition de rester alimenté par la chaîne finale et non par des branchements intermédiaires hérités d’anciennes directions du projet. Une deuxième perspective réside dans l’enrichissement de chambre\_d\_analyse, qui pourrait introduire des mécanismes de corrélation temporelle ou de regroupement narratif d’événements.

Une troisième perspective consisterait à étendre la couverture des tracepoints suivis, mais cette évolution devrait rester compatible avec la philosophie qui a permis la consolidation. Chaque nouvel ajout devrait être justifié à la fois par sa valeur de détection, par son coût de maintenance, et par son impact sur la lisibilité de l’ensemble. L’un des acquis du nettoyage final est justement d’avoir montré que la qualité d’un projet de télémétrie système dépend autant de sa discipline architecturale que de son ambition fonctionnelle.

\appendix

\chapter{Remarques sur l’état final du dépôt}
L’état final du dépôt retient une arborescence francisée cohérente et supprime les doublons qui entretenaient la confusion entre noms anglais et noms français. Cette décision simplifie la lecture, réduit le coût d’onboarding et améliore la fiabilité des usages quotidiens. Elle mérite d’être conservée comme principe pour les évolutions futures, faute de quoi le dépôt pourrait rapidement retomber dans le type d’ambiguïté qui a rendu cette consolidation nécessaire.

\chapter{Remarques sur la reproductibilité}
Les observations présentées dans ce rapport restent reproductibles à deux niveaux. Les essais sur corpus d’exemple peuvent être relancés sans privilège root et servent de garde-fou rapide après une modification du pipeline. Les vérifications positives de capture réelle exigent pour leur part un noyau Linux exposant tracefs ainsi qu’un contexte autorisé à écrire sous le répertoire de traçage. Cette distinction entre validation hors ligne et exécution système doit être comprise comme un trait constitutif du projet et non comme une limitation accidentelle.

\end{document}
